% 3_methods.tex
\section{Methods}
\label{sec:methods}

\subsection{Flux power spectrum}
\label{sec:methods:pk}

For each sightline we clip float-underflow negatives, normalize the flux, apply a
Hann window, take the real FFT along the velocity axis, and square the modulus.
Under the primary per-sightline normalization the flux contrast is
\begin{equation}
  \delta_F = \frac{F}{\langle F\rangle_\mathrm{los}} - 1,
  \label{eq:deltaf}
\end{equation}
where $\langle F\rangle_\mathrm{los}$ is the mean transmitted flux of that
sightline. Dividing by $\langle F\rangle_\mathrm{los}$ removes the per-sightline
mean and forces the classifier onto the \emph{shape} of $P_F(k)$ rather than the
mean-flux offset. The raw positive-frequency power is averaged into
$\nkbins$ log-spaced $k$-bins spanning $k_\mathrm{min}\approx\kmin$ to
$k_\mathrm{Nyq}\approx\knyq~\mathrm{s/km}$, giving a $\nkbins$-dimensional feature
per sightline. We do not apply the $1/\sum w^2$ window correction: the constant
cancels across sightlines and is irrelevant for classification. Per-class
normalization is rejected as label-leaky.

\subsection{Stacking}
\label{sec:methods:stack}

Within each class we average $P_F(k)$ over $M$ disjoint randomly drawn sightlines
to form one stacked spectrum, repeating over the $\Nsightlines/M$ groups. For
independent sightlines the variance of the stacked $P_F(k)$ falls as $1/M$, so
stacking is the signal-to-noise operation appropriate to a power spectrum. We
sweep $M\in\Msweep$. Stacking changes the unit of inference from a single
sightline to a within-class population of $M$-sightline groups; the per-sightline
regime is recovered as the $M=1$ column. This move makes a weak per-sightline
signal detectable as a population statistic, and we frame all stacked results as
population statistics, not single-sightline classification.

\subsection{Dedicated binary classifier}
\label{sec:methods:clf}

For each recipe pair we train a dedicated binary random forest
\citep{Breiman2001} on the stacked $P_F(k)$ features, with $\rfTrees$ trees. We
use $\nfolds$-fold stratified cross-validation, splitting at the post-stacking
group level so that no source sightline appears in both the train and test sets
of a fold. We repeat over $\nseeds$ seeds, where each seed fixes a distinct
stack-membership partition and a distinct fold assignment, giving
$\nseedfold$ balanced-accuracy values per cell. The reported number per cell is
the $16$th percentile (p16) of those $\nseedfold$ values, a conservative
lower-edge summary rather than the mean. Using one dedicated binary classifier
per pair, rather than reading pairwise accuracy off a multiclass confusion
sub-block, gives one estimator and one number per cell, fixed before the run.

The distinguishability lattice is the $4\times4$ table of these p16 values over
the six recipe pairs, plus a strong-AGN-versus-rest detection cell in which
class~4 is classified against classes~1--3 pooled. We adopt a balanced-accuracy
floor of $\latQualBar$ for a pair to count as separable and $\detQualBar$ for the
detection cell, both fixed in advance; these are project-internal operational
gates, not externally calibrated benchmarks.

\subsection{Controls}
\label{sec:methods:controls}

Three controls license the interpretation of a weak signal.

\textbf{Within-class self-pair null.} We split one class into two disjoint halves
of source sightlines, stack each independently, and ask the same classifier to
separate the two halves. The two halves share physics and differ only by
sampling, so any accuracy above chance is a stacking-variance artifact. This
control isolates the contribution of finite-$M$ variance from the recipe signal.

\textbf{Mean-flux-only baseline.} We replace the $\nkbins$-dimensional $P_F(k)$
feature with the single scalar per-class mean flux (and its variance), run the
identical estimator, cross-validation, and stacking, and compare its p16 to the
$P_F(k)$ p16 cell by cell. The headroom $\Delta=p16_{P_F(k)}-p16_{\langle
F\rangle}$ measures how much separability $P_F(k)$ shape carries beyond what mean
flux alone provides.

\textbf{Injection-recovery sensitivity.} We calibrate the smallest $P_F(k)$-shape
perturbation the pipeline can detect at a given $M$. We build a template
$\Delta P(k)=\overline{P}_{C4}(k)-\overline{P}_{C2}(k)$ from the strong-AGN minus
stellar-wind mean spectra, inject $\alpha\,\Delta P(k)$ into per-sightline
stellar-wind spectra before stacking, and classify injected against control
sightlines drawn from disjoint source halves. The minimum detectable amplitude is
\begin{equation}
  \delta_\mathrm{min}(M) = \min\{\alpha : \mathrm{p16}(\alpha) \geq \hfourThresh\},
  \label{eq:deltamin}
\end{equation}
and the $\alpha=0$ arm must center at chance for the instrument to be valid. The
injection is in $P_F(k)$ space, so the resulting sensitivity is strictly a
statement about $P_F(k)$-shape perturbations of amplitude $\alpha\,\Delta P(k)$,
not about feedback physics of strength $\alpha$. We also project the real
stellar-wind/wind${+}$AGN spectral difference onto the injection template,
$\beta=\langle\Delta P_{C3-C2},\Delta P_{C4-C2}\rangle/\lVert\Delta
P_{C4-C2}\rVert^2$, to test whether that difference lies along the
calibrated direction.
